% 图契约：戴维南电压源与诺顿电流源的端口等效；图中的端口、几何、箭头和方程对应本题。
% 版面：沿用前三模块的几何图+推导分区；目标印宽17cm，简图至中等密度。
% 布局：左侧真实结构，右侧关键关系；真实电路边不作为架构流程箭头处理。
\documentclass[tikz,border=7pt]{standalone}
\usepackage{fontspec,xeCJK,amsmath,amssymb}
\IfFontExistsTF{Times New Roman}{\setmainfont{Times New Roman}}{\setmainfont{TeX Gyre Termes}}
\IfFontExistsTF{Heiti SC}{\setCJKmainfont[AutoFakeBold=1.5]{Heiti SC}}{\setCJKmainfont[AutoFakeBold=1.5]{FandolHei-Regular}}
\usetikzlibrary{arrows.meta,bending,calc,patterns,positioning,decorations.pathmorphing}
\definecolor{ink}{HTML}{243745}
\definecolor{blue}{HTML}{275DAB}
\definecolor{green}{HTML}{177C63}
\definecolor{red}{HTML}{BD4935}
\definecolor{orange}{HTML}{AA6C16}
\definecolor{purple}{HTML}{715393}
\tikzset{every node/.style={font=\fontsize{9.3}{12.5}\selectfont,text=ink},
arrow/.style={-{Stealth[length=5pt 1.5,width'=0pt 3]},shorten >=2pt,shorten <=1pt,line width=1pt},
arrow short/.style={-{Stealth[length=3pt 1.0,width'=0pt 2.5]},shorten >=1pt,shorten <=0pt,line width=.8pt},
axis/.style={arrow,draw=ink!60,line width=.6pt},
vec/.style={arrow,draw=blue},
wire/.style={draw=ink!75,line width=.8pt},
guide/.style={dashed,draw=ink!45,line width=.5pt},
card/.style={rounded corners=2pt,draw=blue!65,fill=blue!4,align=center,text width=6.5cm,inner xsep=9pt,inner ysep=7pt},
title/.style={font=\fontsize{10.5}{14}\selectfont\bfseries,text width=15cm,align=center},
label/.style={font=\fontsize{8.4}{11}\selectfont,align=center},
note/.style={label,text width=6.1cm}}
\begin{document}
\begin{tikzpicture}[x=1cm,y=1cm]
\path[use as bounding box] (-8,-3.5) rectangle(8,3.65);
\node[title] at (0,3.25) {\quad 戴维南电压源与诺顿电流源的端口等效};
\begin{scope}[shift={(-4.7,.3)}]\draw[wire,fill=white](-2,0)circle(0.3);\node at(-2,0){};
\node at(-2,.14){$+$};\node at(-2,-.14){$-$};\node[left]at(-2.4,0){$\mathcal E$};\draw[wire](-2,.3)--(-2,1.2)(-2,-.3)--(-2,-1.2)--(2,-1.2);\draw[wire,draw=ink](-2,1.2)--(-0.275,1.2) (0.275,1.2)--(2,1.2);\begin{scope}[shift={(0.0,1.2)},rotate=0.0]\draw[draw=ink,fill=white,line width=.8pt](-0.275,-.085)rectangle(0.275,.085);\end{scope}
\node[label]at(0.0,1.6199999999999999){$r$};
\fill[ink] (2,1.2)circle(2pt);\node[right=4pt] at(2,1.2){$A$};
\fill[ink] (2,-1.2)circle(2pt);\node[right=4pt] at(2,-1.2){$B$};
\draw[<->,orange](2.9,-1.2)--(2.9,1.2)node[midway,right]{$U_{AB}$};\end{scope}\begin{scope}[shift={(3.8,.3)}]\draw[wire,fill=white](-1.3,0)circle(0.3);\node at(-1.3,0){};
\draw[vec,red](-1.3,-.2)--(-1.3,.2);\node[left]at(-1.65,0){$I_N$};\draw[wire](-1.3,.3)--(-1.3,1.2)--(2,1.2)(-1.3,-.3)--(-1.3,-1.2)--(2,-1.2);\draw[wire,draw=ink](0.7,1.2)--(0.7,0.275) (0.7,-0.275)--(0.7,-1.2);\begin{scope}[shift={(0.7,0.0)},rotate=-90.0]\draw[draw=ink,fill=white,line width=.8pt](-0.275,-.085)rectangle(0.275,.085);\end{scope}
\node[label]at(1.12,0.0){$R_N$};
\fill[ink] (0.7,1.2)circle(2pt);
\fill[ink] (0.7,-1.2)circle(2pt);
\fill[ink] (2,1.2)circle(2pt);\node[right=4pt] at(2,1.2){$A$};
\fill[ink] (2,-1.2)circle(2pt);\node[right=4pt] at(2,-1.2){$B$};
\end{scope}\node[card,text width=14cm,draw=blue!65,fill=blue!4]at(0,-2.5){$\mathcal E=U_{\rm oc},\quad I_N=I_{\rm sc}=\mathcal E/r,\quad R_N=R_{\rm th}=r$};

\end{tikzpicture}\end{document}